\section{Elliptic-curve groups and ECDH}

\subsection{From a cubic equation to a group}

For an odd prime $p$, a short Weierstrass curve has the form
\[
  E:\quad y^2=x^3+ax+b,
  \qquad 4a^3+27b^2\not\equiv0\pmod p.
\]
The nonzero discriminant excludes cusps and self-intersections. The solutions
in $\F_p^2$, together with a formal point at infinity $\OO$, form a finite
abelian group. The familiar chord-and-tangent picture over the reals is useful
intuition; over $\F_p$ the curve is a discrete point cloud and the algebraic
formulas are what survive.

Curve25519 uses Montgomery form \cite{montgomery1987}:
\[
  E_{A,B}:\quad Bv^2=u^3+Au^2+u,
  \qquad B(A^2-4)\ne0.
\]
For $B=1$ and distinct noninverse points
$P=(u_1,v_1)$ and $Q=(u_2,v_2)$, put
\[
  \lambda=(v_2-v_1)(u_2-u_1)^{-1}.
\]
Then $P+Q=(u_3,v_3)$ where
\[
  u_3=\lambda^2-A-u_1-u_2,
  \qquad
  v_3=\lambda(u_1-u_3)-v_1.
\]
For doubling, replace the slope by
\[
  \lambda=(3u_1^2+2Au_1+1)(2v_1)^{-1}.
\]
All operations occur in $\F_p$. The exceptional cases complete the group law:
$P+\OO=P$, and $P+(-P)=\OO$ with $-(u,v)=(u,-v)$.

\subsection{Scalar multiplication and ECDLP}

Double-and-add evaluates $[k]P$ from the binary expansion of $k$ using
$O(\log k)$ additions. The corresponding elliptic-curve discrete logarithm
problem is
\[
  (E,G,Q=[a]G)\longmapsto a\pmod{\ord(G)}.
\]
It is exactly the DLP from the preceding section, instantiated in an elliptic
curve subgroup.

Hasse's theorem bounds the ambient group order:
\[
  \bigl|\#E(\F_p)-(p+1)\bigr|\le2\sqrt p.
\]
Cryptographic parameters select a point $G$ of large prime order $\ell$ and
write
\[
  \#E(\F_p)=h\ell,
\]
where $h$ is the cofactor. Multiplying an arbitrary point by $h$ moves it into
a subgroup whose order divides $\ell$, but protocol-specific validation rules
are still required.

\subsection{An exact toy Montgomery curve}

The notebooks use
\[
  E/\F_{127}:\quad v^2=u^3+5u^2+u.
\]
It has $136=8\cdot17$ points. The point
\[
  G=(18,55)
\]
has order $17$, so its subgroup is small enough to enumerate yet has the same
``cofactor times prime'' shape as Curve25519. For example,
\[
  [5]G=(122,73),\qquad [9]G=(124,74),\qquad[11]G=(70,64).
\]
This curve is pedagogical and cryptographically worthless.

\subsection{Elliptic-curve Diffie--Hellman}

Alice chooses $a$ and publishes $A=[a]G$. Bob chooses $b$ and publishes
$B=[b]G$. Associativity of scalar multiplication gives
\[
  [a]B=[a]([b]G)=[ab]G=[b]([a]G)=[b]A.
\]
Alice and Bob never run the hard inverse. They compute \emph{sideways}: each
combines a scalar they know with the public point they received.
\[
\begin{aligned}
  \text{Alice: }(a,B)&\longmapsto[a]B=[ab]G,\\
  \text{Bob: }(b,A)&\longmapsto[b]A=[ab]G,\\
  \text{Eve: }(G,A,B)&\mathrel{\not\!\longrightarrow}[ab]G
    \qquad\text{(the CDH challenge).}
\end{aligned}
\]
The secret scalar is not a trapdoor that makes inversion easy; it enables a
different efficient group action that both honest parties can make agree.

For the toy defaults $a=5$ and $b=9$, both sides reach
\[
  [45]G=[11]G=(70,64),
\]
because scalars are interpreted modulo $17$.

The original Diffie--Hellman insight is that public values can be combined
with a private scalar to reach the same hidden group element
\cite{diffie1976}. A passive observer receives $G,A,B$ and faces CDH; recovering
either scalar via ECDLP is one sufficient way to solve it.
